% Course preamble for "AI on Edge Devices" (KAUST Academy).
%
% Load this file in the main file of every new deck:
%
%   \documentclass[10pt]{beamer}
%   % Course preamble for "AI on Edge Devices" (KAUST Academy).
%
% Load this file in the main file of every new deck:
%
%   \documentclass[10pt]{beamer}
%   % Course preamble for "AI on Edge Devices" (KAUST Academy).
%
% Load this file in the main file of every new deck:
%
%   \documentclass[10pt]{beamer}
%   % Course preamble for "AI on Edge Devices" (KAUST Academy).
%
% Load this file in the main file of every new deck:
%
%   \documentclass[10pt]{beamer}
%   \input{preamble/course}
%
% Do not load preamble/packages.tex or preamble/commands.tex. Those two files
% are the older preamble of the template. No deck uses them, and they report
% LaTeX errors.
%
% This file gives these commands:
%
%   \source{...}       small source line at the bottom of the frame
%   \sourcehere{...}   small source line at the current position
%   \code{...}         inline code
%   \cmark, \xmark     check mark and cross mark, for comparison tables
%   \keypoint{...}     box for the main point of a frame
%   \exerciseframe{title}{body}   one exercise frame
%   \answerframe{title}{body}     the answer frame of that exercise
%   \theorypart{title}{exercise}  start of one part of a theory deck
%   exerciseblock                 environment: the exercise box alone
%   answerblock                   environment: the answer box alone
%   L{width}                      table column: paragraph, ragged right
%   \coursecredits                source list of the credits frame
%
% Listing styles for code: python, cpp, shell. Example:
%
%   \begin{frame}[fragile]{Read the sensor}
%     \begin{lstlisting}[style=python]
%   print("hello")
%     \end{lstlisting}
%   \end{frame}
%
% A frame that contains a listing needs the option [fragile]. A listing also
% cannot stand in the argument of a command. So write such a frame in full and
% use the environment exerciseblock in place of \exerciseframe.
%
% sections/credits.tex prints the credits frame.
%
% Check this preamble after a change:
%   bash tools/build_deck.sh Preamble_Test.tex

% ---------------------------------------------------------------- packages ---
\usepackage[utf8]{inputenc}
\usepackage{amssymb}
\usepackage{amsmath}
\usepackage{graphicx}
\usepackage{caption}
\usepackage{subcaption}
\usepackage{array}
\usepackage{booktabs}
\usepackage{multicol}
\usepackage{multirow}
\usepackage{pifont}
\usepackage{tcolorbox}
\usepackage{listings}
\usepackage{tikz}
\usetikzlibrary{arrows.meta,positioning,fit,calc,shapes.geometric}

% The theme, the logo, and the colours of the template.
\input{preamble/beamer_settings}

% ------------------------------------------------------------ frame layout ---
% Captions on a slide carry no number.
\setbeamertemplate{caption}{\insertcaption}
\setbeamerfont{caption}{size=\footnotesize}
\hypersetup{colorlinks=true,linkcolor=cardinalred,urlcolor=darkcyan}

% The cover frame (sections/cover.tex) has room for a title of one line. A
% longer title moves the date into the footer. This hook counts the lines of
% the title and moves the cover up by the height of each extra line.
\newcount\covertitlelines
\newlength{\covertitleskip}
\addtobeamertemplate{title page}{%
  \setbox0=\vbox{%
    \hsize=\dimexpr\textwidth-16pt\relax
    \usebeamerfont{title}\centering\inserttitle\par
    \global\covertitlelines=\prevgraf
    \global\covertitleskip=\baselineskip}%
  \ifnum\covertitlelines>1
    \vskip-\dimexpr\covertitleskip*\numexpr\covertitlelines-1\relax\relax
  \fi
}{}

% Table column L{width}: a paragraph column with a ragged right margin. A
% plain p{width} column justifies the text and breaks words in a narrow
% column.
\newcolumntype{L}[1]{>{\raggedright\arraybackslash}p{#1}}

% ------------------------------------------------------------ code listings ---
% beamercolorthemestanford.sty gives cardinalred, coolgray, treegreen,
% darkcyan, boxgray, and footergray.
\definecolor{codebg}{RGB}{246,245,243}

\lstdefinestyle{coursebase}{
  basicstyle=\ttfamily\scriptsize,
  keywordstyle=\color{cardinalred}\bfseries,
  commentstyle=\color{coolgray}\itshape,
  stringstyle=\color{darkcyan},
  numberstyle=\ttfamily\tiny\color{coolgray},
  backgroundcolor=\color{codebg},
  frame=single,
  rulecolor=\color{footergray},
  framesep=3pt,
  xleftmargin=2pt,
  xrightmargin=2pt,
  aboveskip=0.6em,
  belowskip=0.4em,
  showstringspaces=false,
  breaklines=true,
  breakatwhitespace=true,
  columns=fullflexible,
  keepspaces=true,
  upquote=true,
  tabsize=2,
}

\lstdefinestyle{python}{
  style=coursebase,
  language=Python,
  morekeywords={as,with,assert,None,True,False,self,async,await},
}

\lstdefinestyle{cpp}{
  style=coursebase,
  language=C++,
  morekeywords={int8_t,uint8_t,int16_t,uint16_t,int32_t,uint32_t,size_t,
                constexpr,nullptr,alignas,PROGMEM},
}

\lstdefinestyle{shell}{
  style=coursebase,
  language=sh,
  morekeywords={python3,pip,pip3,arduino-cli,mpremote,esptool.py,ollama,
                mosquitto_pub,mosquitto_sub,ffmpeg,systemctl,raspi-config},
}

% Plain \begin{lstlisting} uses the base style.
\lstset{style=coursebase}

% ---------------------------------------------------------------- commands ---
% Inline code.
\newcommand{\code}[1]{\texttt{\small #1}}

% Check mark and cross mark.
\newcommand{\cmark}{{\color{treegreen}\ding{51}}}
\newcommand{\xmark}{{\color{cardinalred}\ding{55}}}

% Source line for a figure, a table, or text from a source.
% \sourcehere prints the line at the current position, for example under a
% figure in one column. \source pushes the line to the bottom of the frame.
\newcommand{\sourcehere}[1]{{\par\scriptsize\color{coolgray}Source: #1\par}}
\newcommand{\source}[1]{\vfill\sourcehere{#1}}

% The main point of a frame.
\newcommand{\keypoint}[1]{%
  \begin{tcolorbox}[colframe=cardinalred,colback=boxgray,boxrule=0.6pt,
      arc=2pt,left=4pt,right=4pt,top=3pt,bottom=3pt]
    \textbf{Key point:} #1
  \end{tcolorbox}%
}

% ------------------------------------------------------ exercise and answer ---
% \exerciseframe and \answerframe make the complete frame. They read the body
% as a macro argument, so the body must contain no listing.
% The environments exerciseblock and answerblock hold the content alone. Use
% them inside a frame with the option [fragile] when the body has a listing.
\newtcolorbox{exerciseblock}{colframe=darkcyan,colback=boxgray,boxrule=0.8pt,
  arc=2pt,left=5pt,right=5pt,top=4pt,bottom=4pt}

\newtcolorbox{answerblock}{colframe=treegreen,colback=boxgray,boxrule=0.8pt,
  arc=2pt,left=5pt,right=5pt,top=4pt,bottom=4pt}

\newcommand{\exerciseframe}[2]{%
  \begin{frame}{Exercise: #1}
    \begin{exerciseblock}#2\end{exerciseblock}
  \end{frame}%
}

\newcommand{\answerframe}[2]{%
  \begin{frame}{Answer: #1}
    \begin{answerblock}#2\end{answerblock}
  \end{frame}%
}

% ------------------------------------------------------------ theory parts ---
% A theory deck has three parts. Each part file starts with
%   \theorypart{TITLE}{EXERCISE}
% The command starts a section of the deck. The contents frame
% (sections/dayNN/toc.tex) lists the title and the exercise of each part.
% Before part 2 and before part 3 the contents frame comes again, with the
% current part highlighted. Part 1 has no such frame, because the contents
% frame of the deck stands immediately before it.
\DeclareRobustCommand{\partexercise}[1]{\\{\footnotesize Exercise: #1}}
\newcommand{\theorypart}[2]{%
  \section[#1]{#1\texorpdfstring{\partexercise{#2}}{}}}

\setbeamertemplate{section in toc}{%
  \makebox[3.8em][l]{\textbf{Part \inserttocsectionnumber}}%
  \parbox[t]{\dimexpr\linewidth-3.8em\relax}{\raggedright\inserttocsection}}
\setbeamertemplate{section in toc shaded}[default][45]

\AtBeginSection{%
  \ifnum\value{section}>1
    \begin{frame}{Contents}
      \tableofcontents[currentsection]
    \end{frame}
  \fi}

% ----------------------------------------------------------------- credits ---
% One command for each source. Section 10 of the syllabus gives this text.
\newcommand{\creditmlsys}{\item \emph{Machine Learning Systems} by Vijay
  Janapa Reddi and contributors, Harvard University. mlsysbook.ai,
  CC BY-NC-SA 4.0.}
\newcommand{\creditraspi}{\item \emph{Edge AI Engineering: Raspberry Pi} by
  Marcelo Rovai. github.com/Mjrovai/EdgeML\_Made\_Ease\_ebook and
  github.com/Mjrovai/EdgeML-with-Raspberry-Pi, GPL-3.0.}
\newcommand{\creditxiao}{\item \emph{TinyML Made Easy: XIAO ESP32S3} by
  Marcelo Rovai. github.com/Mjrovai/XIAO-ESP32S3-Sense, Apache-2.0.}
\newcommand{\creditbigpower}{\item \emph{XIAO: Big Power, Small Board} by Lei
  Feng and Marcelo Rovai.
  github.com/Mjrovai/XIAO\_Big\_Power\_Small\_Board-ebook, GPL-3.0.}
\newcommand{\credittinymlx}{\item \emph{HarvardX TinyML courseware} by the
  TinyMLx team. github.com/tinyMLx/courseware, CC BY-NC-SA 4.0.}
\newcommand{\credittflm}{\item \emph{TensorFlow Lite Micro Arduino examples}
  by The TensorFlow Authors.
  github.com/tensorflow/tflite-micro-arduino-examples, Apache-2.0.}

% The source list of the credits frame. The default names the main textbook
% and the three companion books. A deck changes the list, so that the frame
% names only the sources that the deck uses:
%   \renewcommand{\coursecredits}{\creditmlsys\creditxiao}
\newcommand{\coursecredits}{%
  \creditmlsys
  \creditraspi
  \creditxiao
  \creditbigpower
}

%
% Do not load preamble/packages.tex or preamble/commands.tex. Those two files
% are the older preamble of the template. No deck uses them, and they report
% LaTeX errors.
%
% This file gives these commands:
%
%   \source{...}       small source line at the bottom of the frame
%   \sourcehere{...}   small source line at the current position
%   \code{...}         inline code
%   \cmark, \xmark     check mark and cross mark, for comparison tables
%   \keypoint{...}     box for the main point of a frame
%   \exerciseframe{title}{body}   one exercise frame
%   \answerframe{title}{body}     the answer frame of that exercise
%   \theorypart{title}{exercise}  start of one part of a theory deck
%   exerciseblock                 environment: the exercise box alone
%   answerblock                   environment: the answer box alone
%   L{width}                      table column: paragraph, ragged right
%   \coursecredits                source list of the credits frame
%
% Listing styles for code: python, cpp, shell. Example:
%
%   \begin{frame}[fragile]{Read the sensor}
%     \begin{lstlisting}[style=python]
%   print("hello")
%     \end{lstlisting}
%   \end{frame}
%
% A frame that contains a listing needs the option [fragile]. A listing also
% cannot stand in the argument of a command. So write such a frame in full and
% use the environment exerciseblock in place of \exerciseframe.
%
% sections/credits.tex prints the credits frame.
%
% Check this preamble after a change:
%   bash tools/build_deck.sh Preamble_Test.tex

% ---------------------------------------------------------------- packages ---
\usepackage[utf8]{inputenc}
\usepackage{amssymb}
\usepackage{amsmath}
\usepackage{graphicx}
\usepackage{caption}
\usepackage{subcaption}
\usepackage{array}
\usepackage{booktabs}
\usepackage{multicol}
\usepackage{multirow}
\usepackage{pifont}
\usepackage{tcolorbox}
\usepackage{listings}
\usepackage{tikz}
\usetikzlibrary{arrows.meta,positioning,fit,calc,shapes.geometric}

% The theme, the logo, and the colours of the template.
% Beamer settings
\usetheme{Stanford} 
\input{style_files/my_beamer_defs.sty}
\logo{
    \includegraphics[height=0.4in]{style_files/kaust_academy_logo.png}
    % \includegraphics[height=0.4in]{style_files/LMHLogo.jpeg}
}

% Commands to relax beamer and subfig conflicts
\makeatletter
\let\@@magyar@captionfix\relax
\makeatother

% Set itemize styles
\setbeamertemplate{itemize items}[default]
\setbeamertemplate{itemize subitem}[circle] 

% ------------------------------------------------------------ frame layout ---
% Captions on a slide carry no number.
\setbeamertemplate{caption}{\insertcaption}
\setbeamerfont{caption}{size=\footnotesize}
\hypersetup{colorlinks=true,linkcolor=cardinalred,urlcolor=darkcyan}

% The cover frame (sections/cover.tex) has room for a title of one line. A
% longer title moves the date into the footer. This hook counts the lines of
% the title and moves the cover up by the height of each extra line.
\newcount\covertitlelines
\newlength{\covertitleskip}
\addtobeamertemplate{title page}{%
  \setbox0=\vbox{%
    \hsize=\dimexpr\textwidth-16pt\relax
    \usebeamerfont{title}\centering\inserttitle\par
    \global\covertitlelines=\prevgraf
    \global\covertitleskip=\baselineskip}%
  \ifnum\covertitlelines>1
    \vskip-\dimexpr\covertitleskip*\numexpr\covertitlelines-1\relax\relax
  \fi
}{}

% Table column L{width}: a paragraph column with a ragged right margin. A
% plain p{width} column justifies the text and breaks words in a narrow
% column.
\newcolumntype{L}[1]{>{\raggedright\arraybackslash}p{#1}}

% ------------------------------------------------------------ code listings ---
% beamercolorthemestanford.sty gives cardinalred, coolgray, treegreen,
% darkcyan, boxgray, and footergray.
\definecolor{codebg}{RGB}{246,245,243}

\lstdefinestyle{coursebase}{
  basicstyle=\ttfamily\scriptsize,
  keywordstyle=\color{cardinalred}\bfseries,
  commentstyle=\color{coolgray}\itshape,
  stringstyle=\color{darkcyan},
  numberstyle=\ttfamily\tiny\color{coolgray},
  backgroundcolor=\color{codebg},
  frame=single,
  rulecolor=\color{footergray},
  framesep=3pt,
  xleftmargin=2pt,
  xrightmargin=2pt,
  aboveskip=0.6em,
  belowskip=0.4em,
  showstringspaces=false,
  breaklines=true,
  breakatwhitespace=true,
  columns=fullflexible,
  keepspaces=true,
  upquote=true,
  tabsize=2,
}

\lstdefinestyle{python}{
  style=coursebase,
  language=Python,
  morekeywords={as,with,assert,None,True,False,self,async,await},
}

\lstdefinestyle{cpp}{
  style=coursebase,
  language=C++,
  morekeywords={int8_t,uint8_t,int16_t,uint16_t,int32_t,uint32_t,size_t,
                constexpr,nullptr,alignas,PROGMEM},
}

\lstdefinestyle{shell}{
  style=coursebase,
  language=sh,
  morekeywords={python3,pip,pip3,arduino-cli,mpremote,esptool.py,ollama,
                mosquitto_pub,mosquitto_sub,ffmpeg,systemctl,raspi-config},
}

% Plain \begin{lstlisting} uses the base style.
\lstset{style=coursebase}

% ---------------------------------------------------------------- commands ---
% Inline code.
\newcommand{\code}[1]{\texttt{\small #1}}

% Check mark and cross mark.
\newcommand{\cmark}{{\color{treegreen}\ding{51}}}
\newcommand{\xmark}{{\color{cardinalred}\ding{55}}}

% Source line for a figure, a table, or text from a source.
% \sourcehere prints the line at the current position, for example under a
% figure in one column. \source pushes the line to the bottom of the frame.
\newcommand{\sourcehere}[1]{{\par\scriptsize\color{coolgray}Source: #1\par}}
\newcommand{\source}[1]{\vfill\sourcehere{#1}}

% The main point of a frame.
\newcommand{\keypoint}[1]{%
  \begin{tcolorbox}[colframe=cardinalred,colback=boxgray,boxrule=0.6pt,
      arc=2pt,left=4pt,right=4pt,top=3pt,bottom=3pt]
    \textbf{Key point:} #1
  \end{tcolorbox}%
}

% ------------------------------------------------------ exercise and answer ---
% \exerciseframe and \answerframe make the complete frame. They read the body
% as a macro argument, so the body must contain no listing.
% The environments exerciseblock and answerblock hold the content alone. Use
% them inside a frame with the option [fragile] when the body has a listing.
\newtcolorbox{exerciseblock}{colframe=darkcyan,colback=boxgray,boxrule=0.8pt,
  arc=2pt,left=5pt,right=5pt,top=4pt,bottom=4pt}

\newtcolorbox{answerblock}{colframe=treegreen,colback=boxgray,boxrule=0.8pt,
  arc=2pt,left=5pt,right=5pt,top=4pt,bottom=4pt}

\newcommand{\exerciseframe}[2]{%
  \begin{frame}{Exercise: #1}
    \begin{exerciseblock}#2\end{exerciseblock}
  \end{frame}%
}

\newcommand{\answerframe}[2]{%
  \begin{frame}{Answer: #1}
    \begin{answerblock}#2\end{answerblock}
  \end{frame}%
}

% ------------------------------------------------------------ theory parts ---
% A theory deck has three parts. Each part file starts with
%   \theorypart{TITLE}{EXERCISE}
% The command starts a section of the deck. The contents frame
% (sections/dayNN/toc.tex) lists the title and the exercise of each part.
% Before part 2 and before part 3 the contents frame comes again, with the
% current part highlighted. Part 1 has no such frame, because the contents
% frame of the deck stands immediately before it.
\DeclareRobustCommand{\partexercise}[1]{\\{\footnotesize Exercise: #1}}
\newcommand{\theorypart}[2]{%
  \section[#1]{#1\texorpdfstring{\partexercise{#2}}{}}}

\setbeamertemplate{section in toc}{%
  \makebox[3.8em][l]{\textbf{Part \inserttocsectionnumber}}%
  \parbox[t]{\dimexpr\linewidth-3.8em\relax}{\raggedright\inserttocsection}}
\setbeamertemplate{section in toc shaded}[default][45]

\AtBeginSection{%
  \ifnum\value{section}>1
    \begin{frame}{Contents}
      \tableofcontents[currentsection]
    \end{frame}
  \fi}

% ----------------------------------------------------------------- credits ---
% One command for each source. Section 10 of the syllabus gives this text.
\newcommand{\creditmlsys}{\item \emph{Machine Learning Systems} by Vijay
  Janapa Reddi and contributors, Harvard University. mlsysbook.ai,
  CC BY-NC-SA 4.0.}
\newcommand{\creditraspi}{\item \emph{Edge AI Engineering: Raspberry Pi} by
  Marcelo Rovai. github.com/Mjrovai/EdgeML\_Made\_Ease\_ebook and
  github.com/Mjrovai/EdgeML-with-Raspberry-Pi, GPL-3.0.}
\newcommand{\creditxiao}{\item \emph{TinyML Made Easy: XIAO ESP32S3} by
  Marcelo Rovai. github.com/Mjrovai/XIAO-ESP32S3-Sense, Apache-2.0.}
\newcommand{\creditbigpower}{\item \emph{XIAO: Big Power, Small Board} by Lei
  Feng and Marcelo Rovai.
  github.com/Mjrovai/XIAO\_Big\_Power\_Small\_Board-ebook, GPL-3.0.}
\newcommand{\credittinymlx}{\item \emph{HarvardX TinyML courseware} by the
  TinyMLx team. github.com/tinyMLx/courseware, CC BY-NC-SA 4.0.}
\newcommand{\credittflm}{\item \emph{TensorFlow Lite Micro Arduino examples}
  by The TensorFlow Authors.
  github.com/tensorflow/tflite-micro-arduino-examples, Apache-2.0.}

% The source list of the credits frame. The default names the main textbook
% and the three companion books. A deck changes the list, so that the frame
% names only the sources that the deck uses:
%   \renewcommand{\coursecredits}{\creditmlsys\creditxiao}
\newcommand{\coursecredits}{%
  \creditmlsys
  \creditraspi
  \creditxiao
  \creditbigpower
}

%
% Do not load preamble/packages.tex or preamble/commands.tex. Those two files
% are the older preamble of the template. No deck uses them, and they report
% LaTeX errors.
%
% This file gives these commands:
%
%   \source{...}       small source line at the bottom of the frame
%   \sourcehere{...}   small source line at the current position
%   \code{...}         inline code
%   \cmark, \xmark     check mark and cross mark, for comparison tables
%   \keypoint{...}     box for the main point of a frame
%   \exerciseframe{title}{body}   one exercise frame
%   \answerframe{title}{body}     the answer frame of that exercise
%   \theorypart{title}{exercise}  start of one part of a theory deck
%   exerciseblock                 environment: the exercise box alone
%   answerblock                   environment: the answer box alone
%   L{width}                      table column: paragraph, ragged right
%   \coursecredits                source list of the credits frame
%
% Listing styles for code: python, cpp, shell. Example:
%
%   \begin{frame}[fragile]{Read the sensor}
%     \begin{lstlisting}[style=python]
%   print("hello")
%     \end{lstlisting}
%   \end{frame}
%
% A frame that contains a listing needs the option [fragile]. A listing also
% cannot stand in the argument of a command. So write such a frame in full and
% use the environment exerciseblock in place of \exerciseframe.
%
% sections/credits.tex prints the credits frame.
%
% Check this preamble after a change:
%   bash tools/build_deck.sh Preamble_Test.tex

% ---------------------------------------------------------------- packages ---
\usepackage[utf8]{inputenc}
\usepackage{amssymb}
\usepackage{amsmath}
\usepackage{graphicx}
\usepackage{caption}
\usepackage{subcaption}
\usepackage{array}
\usepackage{booktabs}
\usepackage{multicol}
\usepackage{multirow}
\usepackage{pifont}
\usepackage{tcolorbox}
\usepackage{listings}
\usepackage{tikz}
\usetikzlibrary{arrows.meta,positioning,fit,calc,shapes.geometric}

% The theme, the logo, and the colours of the template.
% Beamer settings
\usetheme{Stanford} 
\input{style_files/my_beamer_defs.sty}
\logo{
    \includegraphics[height=0.4in]{style_files/kaust_academy_logo.png}
    % \includegraphics[height=0.4in]{style_files/LMHLogo.jpeg}
}

% Commands to relax beamer and subfig conflicts
\makeatletter
\let\@@magyar@captionfix\relax
\makeatother

% Set itemize styles
\setbeamertemplate{itemize items}[default]
\setbeamertemplate{itemize subitem}[circle] 

% ------------------------------------------------------------ frame layout ---
% Captions on a slide carry no number.
\setbeamertemplate{caption}{\insertcaption}
\setbeamerfont{caption}{size=\footnotesize}
\hypersetup{colorlinks=true,linkcolor=cardinalred,urlcolor=darkcyan}

% The cover frame (sections/cover.tex) has room for a title of one line. A
% longer title moves the date into the footer. This hook counts the lines of
% the title and moves the cover up by the height of each extra line.
\newcount\covertitlelines
\newlength{\covertitleskip}
\addtobeamertemplate{title page}{%
  \setbox0=\vbox{%
    \hsize=\dimexpr\textwidth-16pt\relax
    \usebeamerfont{title}\centering\inserttitle\par
    \global\covertitlelines=\prevgraf
    \global\covertitleskip=\baselineskip}%
  \ifnum\covertitlelines>1
    \vskip-\dimexpr\covertitleskip*\numexpr\covertitlelines-1\relax\relax
  \fi
}{}

% Table column L{width}: a paragraph column with a ragged right margin. A
% plain p{width} column justifies the text and breaks words in a narrow
% column.
\newcolumntype{L}[1]{>{\raggedright\arraybackslash}p{#1}}

% ------------------------------------------------------------ code listings ---
% beamercolorthemestanford.sty gives cardinalred, coolgray, treegreen,
% darkcyan, boxgray, and footergray.
\definecolor{codebg}{RGB}{246,245,243}

\lstdefinestyle{coursebase}{
  basicstyle=\ttfamily\scriptsize,
  keywordstyle=\color{cardinalred}\bfseries,
  commentstyle=\color{coolgray}\itshape,
  stringstyle=\color{darkcyan},
  numberstyle=\ttfamily\tiny\color{coolgray},
  backgroundcolor=\color{codebg},
  frame=single,
  rulecolor=\color{footergray},
  framesep=3pt,
  xleftmargin=2pt,
  xrightmargin=2pt,
  aboveskip=0.6em,
  belowskip=0.4em,
  showstringspaces=false,
  breaklines=true,
  breakatwhitespace=true,
  columns=fullflexible,
  keepspaces=true,
  upquote=true,
  tabsize=2,
}

\lstdefinestyle{python}{
  style=coursebase,
  language=Python,
  morekeywords={as,with,assert,None,True,False,self,async,await},
}

\lstdefinestyle{cpp}{
  style=coursebase,
  language=C++,
  morekeywords={int8_t,uint8_t,int16_t,uint16_t,int32_t,uint32_t,size_t,
                constexpr,nullptr,alignas,PROGMEM},
}

\lstdefinestyle{shell}{
  style=coursebase,
  language=sh,
  morekeywords={python3,pip,pip3,arduino-cli,mpremote,esptool.py,ollama,
                mosquitto_pub,mosquitto_sub,ffmpeg,systemctl,raspi-config},
}

% Plain \begin{lstlisting} uses the base style.
\lstset{style=coursebase}

% ---------------------------------------------------------------- commands ---
% Inline code.
\newcommand{\code}[1]{\texttt{\small #1}}

% Check mark and cross mark.
\newcommand{\cmark}{{\color{treegreen}\ding{51}}}
\newcommand{\xmark}{{\color{cardinalred}\ding{55}}}

% Source line for a figure, a table, or text from a source.
% \sourcehere prints the line at the current position, for example under a
% figure in one column. \source pushes the line to the bottom of the frame.
\newcommand{\sourcehere}[1]{{\par\scriptsize\color{coolgray}Source: #1\par}}
\newcommand{\source}[1]{\vfill\sourcehere{#1}}

% The main point of a frame.
\newcommand{\keypoint}[1]{%
  \begin{tcolorbox}[colframe=cardinalred,colback=boxgray,boxrule=0.6pt,
      arc=2pt,left=4pt,right=4pt,top=3pt,bottom=3pt]
    \textbf{Key point:} #1
  \end{tcolorbox}%
}

% ------------------------------------------------------ exercise and answer ---
% \exerciseframe and \answerframe make the complete frame. They read the body
% as a macro argument, so the body must contain no listing.
% The environments exerciseblock and answerblock hold the content alone. Use
% them inside a frame with the option [fragile] when the body has a listing.
\newtcolorbox{exerciseblock}{colframe=darkcyan,colback=boxgray,boxrule=0.8pt,
  arc=2pt,left=5pt,right=5pt,top=4pt,bottom=4pt}

\newtcolorbox{answerblock}{colframe=treegreen,colback=boxgray,boxrule=0.8pt,
  arc=2pt,left=5pt,right=5pt,top=4pt,bottom=4pt}

\newcommand{\exerciseframe}[2]{%
  \begin{frame}{Exercise: #1}
    \begin{exerciseblock}#2\end{exerciseblock}
  \end{frame}%
}

\newcommand{\answerframe}[2]{%
  \begin{frame}{Answer: #1}
    \begin{answerblock}#2\end{answerblock}
  \end{frame}%
}

% ------------------------------------------------------------ theory parts ---
% A theory deck has three parts. Each part file starts with
%   \theorypart{TITLE}{EXERCISE}
% The command starts a section of the deck. The contents frame
% (sections/dayNN/toc.tex) lists the title and the exercise of each part.
% Before part 2 and before part 3 the contents frame comes again, with the
% current part highlighted. Part 1 has no such frame, because the contents
% frame of the deck stands immediately before it.
\DeclareRobustCommand{\partexercise}[1]{\\{\footnotesize Exercise: #1}}
\newcommand{\theorypart}[2]{%
  \section[#1]{#1\texorpdfstring{\partexercise{#2}}{}}}

\setbeamertemplate{section in toc}{%
  \makebox[3.8em][l]{\textbf{Part \inserttocsectionnumber}}%
  \parbox[t]{\dimexpr\linewidth-3.8em\relax}{\raggedright\inserttocsection}}
\setbeamertemplate{section in toc shaded}[default][45]

\AtBeginSection{%
  \ifnum\value{section}>1
    \begin{frame}{Contents}
      \tableofcontents[currentsection]
    \end{frame}
  \fi}

% ----------------------------------------------------------------- credits ---
% One command for each source. Section 10 of the syllabus gives this text.
\newcommand{\creditmlsys}{\item \emph{Machine Learning Systems} by Vijay
  Janapa Reddi and contributors, Harvard University. mlsysbook.ai,
  CC BY-NC-SA 4.0.}
\newcommand{\creditraspi}{\item \emph{Edge AI Engineering: Raspberry Pi} by
  Marcelo Rovai. github.com/Mjrovai/EdgeML\_Made\_Ease\_ebook and
  github.com/Mjrovai/EdgeML-with-Raspberry-Pi, GPL-3.0.}
\newcommand{\creditxiao}{\item \emph{TinyML Made Easy: XIAO ESP32S3} by
  Marcelo Rovai. github.com/Mjrovai/XIAO-ESP32S3-Sense, Apache-2.0.}
\newcommand{\creditbigpower}{\item \emph{XIAO: Big Power, Small Board} by Lei
  Feng and Marcelo Rovai.
  github.com/Mjrovai/XIAO\_Big\_Power\_Small\_Board-ebook, GPL-3.0.}
\newcommand{\credittinymlx}{\item \emph{HarvardX TinyML courseware} by the
  TinyMLx team. github.com/tinyMLx/courseware, CC BY-NC-SA 4.0.}
\newcommand{\credittflm}{\item \emph{TensorFlow Lite Micro Arduino examples}
  by The TensorFlow Authors.
  github.com/tensorflow/tflite-micro-arduino-examples, Apache-2.0.}

% The source list of the credits frame. The default names the main textbook
% and the three companion books. A deck changes the list, so that the frame
% names only the sources that the deck uses:
%   \renewcommand{\coursecredits}{\creditmlsys\creditxiao}
\newcommand{\coursecredits}{%
  \creditmlsys
  \creditraspi
  \creditxiao
  \creditbigpower
}

%
% Do not load preamble/packages.tex or preamble/commands.tex. Those two files
% are the older preamble of the template. No deck uses them, and they report
% LaTeX errors.
%
% This file gives these commands:
%
%   \source{...}       small source line at the bottom of the frame
%   \sourcehere{...}   small source line at the current position
%   \code{...}         inline code
%   \cmark, \xmark     check mark and cross mark, for comparison tables
%   \keypoint{...}     box for the main point of a frame
%   \exerciseframe{title}{body}   one exercise frame
%   \answerframe{title}{body}     the answer frame of that exercise
%   \theorypart{title}{exercise}  start of one part of a theory deck
%   exerciseblock                 environment: the exercise box alone
%   answerblock                   environment: the answer box alone
%   L{width}                      table column: paragraph, ragged right
%   \coursecredits                source list of the credits frame
%
% Listing styles for code: python, cpp, shell. Example:
%
%   \begin{frame}[fragile]{Read the sensor}
%     \begin{lstlisting}[style=python]
%   print("hello")
%     \end{lstlisting}
%   \end{frame}
%
% A frame that contains a listing needs the option [fragile]. A listing also
% cannot stand in the argument of a command. So write such a frame in full and
% use the environment exerciseblock in place of \exerciseframe.
%
% sections/credits.tex prints the credits frame.
%
% Check this preamble after a change:
%   bash tools/build_deck.sh Preamble_Test.tex

% ---------------------------------------------------------------- packages ---
\usepackage[utf8]{inputenc}
\usepackage{amssymb}
\usepackage{amsmath}
\usepackage{graphicx}
\usepackage{caption}
\usepackage{subcaption}
\usepackage{array}
\usepackage{booktabs}
\usepackage{multicol}
\usepackage{multirow}
\usepackage{pifont}
\usepackage{tcolorbox}
\usepackage{listings}
\usepackage{tikz}
\usetikzlibrary{arrows.meta,positioning,fit,calc,shapes.geometric}

% The theme, the logo, and the colours of the template.
% Beamer settings
\usetheme{Stanford} 
\input{style_files/my_beamer_defs.sty}
\logo{
    \includegraphics[height=0.4in]{style_files/kaust_academy_logo.png}
    % \includegraphics[height=0.4in]{style_files/LMHLogo.jpeg}
}

% Commands to relax beamer and subfig conflicts
\makeatletter
\let\@@magyar@captionfix\relax
\makeatother

% Set itemize styles
\setbeamertemplate{itemize items}[default]
\setbeamertemplate{itemize subitem}[circle] 

% ------------------------------------------------------------ frame layout ---
% Captions on a slide carry no number.
\setbeamertemplate{caption}{\insertcaption}
\setbeamerfont{caption}{size=\footnotesize}
\hypersetup{colorlinks=true,linkcolor=cardinalred,urlcolor=darkcyan}

% The cover frame (sections/cover.tex) has room for a title of one line. A
% longer title moves the date into the footer. This hook counts the lines of
% the title and moves the cover up by the height of each extra line.
\newcount\covertitlelines
\newlength{\covertitleskip}
\addtobeamertemplate{title page}{%
  \setbox0=\vbox{%
    \hsize=\dimexpr\textwidth-16pt\relax
    \usebeamerfont{title}\centering\inserttitle\par
    \global\covertitlelines=\prevgraf
    \global\covertitleskip=\baselineskip}%
  \ifnum\covertitlelines>1
    \vskip-\dimexpr\covertitleskip*\numexpr\covertitlelines-1\relax\relax
  \fi
}{}

% Table column L{width}: a paragraph column with a ragged right margin. A
% plain p{width} column justifies the text and breaks words in a narrow
% column.
\newcolumntype{L}[1]{>{\raggedright\arraybackslash}p{#1}}

% ------------------------------------------------------------ code listings ---
% beamercolorthemestanford.sty gives cardinalred, coolgray, treegreen,
% darkcyan, boxgray, and footergray.
\definecolor{codebg}{RGB}{246,245,243}

\lstdefinestyle{coursebase}{
  basicstyle=\ttfamily\scriptsize,
  keywordstyle=\color{cardinalred}\bfseries,
  commentstyle=\color{coolgray}\itshape,
  stringstyle=\color{darkcyan},
  numberstyle=\ttfamily\tiny\color{coolgray},
  backgroundcolor=\color{codebg},
  frame=single,
  rulecolor=\color{footergray},
  framesep=3pt,
  xleftmargin=2pt,
  xrightmargin=2pt,
  aboveskip=0.6em,
  belowskip=0.4em,
  showstringspaces=false,
  breaklines=true,
  breakatwhitespace=true,
  columns=fullflexible,
  keepspaces=true,
  upquote=true,
  tabsize=2,
}

\lstdefinestyle{python}{
  style=coursebase,
  language=Python,
  morekeywords={as,with,assert,None,True,False,self,async,await},
}

\lstdefinestyle{cpp}{
  style=coursebase,
  language=C++,
  morekeywords={int8_t,uint8_t,int16_t,uint16_t,int32_t,uint32_t,size_t,
                constexpr,nullptr,alignas,PROGMEM},
}

\lstdefinestyle{shell}{
  style=coursebase,
  language=sh,
  morekeywords={python3,pip,pip3,arduino-cli,mpremote,esptool.py,ollama,
                mosquitto_pub,mosquitto_sub,ffmpeg,systemctl,raspi-config},
}

% Plain \begin{lstlisting} uses the base style.
\lstset{style=coursebase}

% ---------------------------------------------------------------- commands ---
% Inline code.
\newcommand{\code}[1]{\texttt{\small #1}}

% Check mark and cross mark.
\newcommand{\cmark}{{\color{treegreen}\ding{51}}}
\newcommand{\xmark}{{\color{cardinalred}\ding{55}}}

% Source line for a figure, a table, or text from a source.
% \sourcehere prints the line at the current position, for example under a
% figure in one column. \source pushes the line to the bottom of the frame.
\newcommand{\sourcehere}[1]{{\par\scriptsize\color{coolgray}Source: #1\par}}
\newcommand{\source}[1]{\vfill\sourcehere{#1}}

% The main point of a frame.
\newcommand{\keypoint}[1]{%
  \begin{tcolorbox}[colframe=cardinalred,colback=boxgray,boxrule=0.6pt,
      arc=2pt,left=4pt,right=4pt,top=3pt,bottom=3pt]
    \textbf{Key point:} #1
  \end{tcolorbox}%
}

% ------------------------------------------------------ exercise and answer ---
% \exerciseframe and \answerframe make the complete frame. They read the body
% as a macro argument, so the body must contain no listing.
% The environments exerciseblock and answerblock hold the content alone. Use
% them inside a frame with the option [fragile] when the body has a listing.
\newtcolorbox{exerciseblock}{colframe=darkcyan,colback=boxgray,boxrule=0.8pt,
  arc=2pt,left=5pt,right=5pt,top=4pt,bottom=4pt}

\newtcolorbox{answerblock}{colframe=treegreen,colback=boxgray,boxrule=0.8pt,
  arc=2pt,left=5pt,right=5pt,top=4pt,bottom=4pt}

\newcommand{\exerciseframe}[2]{%
  \begin{frame}{Exercise: #1}
    \begin{exerciseblock}#2\end{exerciseblock}
  \end{frame}%
}

\newcommand{\answerframe}[2]{%
  \begin{frame}{Answer: #1}
    \begin{answerblock}#2\end{answerblock}
  \end{frame}%
}

% ------------------------------------------------------------ theory parts ---
% A theory deck has three parts. Each part file starts with
%   \theorypart{TITLE}{EXERCISE}
% The command starts a section of the deck. The contents frame
% (sections/dayNN/toc.tex) lists the title and the exercise of each part.
% Before part 2 and before part 3 the contents frame comes again, with the
% current part highlighted. Part 1 has no such frame, because the contents
% frame of the deck stands immediately before it.
\DeclareRobustCommand{\partexercise}[1]{\\{\footnotesize Exercise: #1}}
\newcommand{\theorypart}[2]{%
  \section[#1]{#1\texorpdfstring{\partexercise{#2}}{}}}

\setbeamertemplate{section in toc}{%
  \makebox[3.8em][l]{\textbf{Part \inserttocsectionnumber}}%
  \parbox[t]{\dimexpr\linewidth-3.8em\relax}{\raggedright\inserttocsection}}
\setbeamertemplate{section in toc shaded}[default][45]

\AtBeginSection{%
  \ifnum\value{section}>1
    \begin{frame}{Contents}
      \tableofcontents[currentsection]
    \end{frame}
  \fi}

% ----------------------------------------------------------------- credits ---
% One command for each source. Section 10 of the syllabus gives this text.
\newcommand{\creditmlsys}{\item \emph{Machine Learning Systems} by Vijay
  Janapa Reddi and contributors, Harvard University. mlsysbook.ai,
  CC BY-NC-SA 4.0.}
\newcommand{\creditraspi}{\item \emph{Edge AI Engineering: Raspberry Pi} by
  Marcelo Rovai. github.com/Mjrovai/EdgeML\_Made\_Ease\_ebook and
  github.com/Mjrovai/EdgeML-with-Raspberry-Pi, GPL-3.0.}
\newcommand{\creditxiao}{\item \emph{TinyML Made Easy: XIAO ESP32S3} by
  Marcelo Rovai. github.com/Mjrovai/XIAO-ESP32S3-Sense, Apache-2.0.}
\newcommand{\creditbigpower}{\item \emph{XIAO: Big Power, Small Board} by Lei
  Feng and Marcelo Rovai.
  github.com/Mjrovai/XIAO\_Big\_Power\_Small\_Board-ebook, GPL-3.0.}
\newcommand{\credittinymlx}{\item \emph{HarvardX TinyML courseware} by the
  TinyMLx team. github.com/tinyMLx/courseware, CC BY-NC-SA 4.0.}
\newcommand{\credittflm}{\item \emph{TensorFlow Lite Micro Arduino examples}
  by The TensorFlow Authors.
  github.com/tensorflow/tflite-micro-arduino-examples, Apache-2.0.}

% The source list of the credits frame. The default names the main textbook
% and the three companion books. A deck changes the list, so that the frame
% names only the sources that the deck uses:
%   \renewcommand{\coursecredits}{\creditmlsys\creditxiao}
\newcommand{\coursecredits}{%
  \creditmlsys
  \creditraspi
  \creditxiao
  \creditbigpower
}
